\section{Model Strengths and Limitations}
\subsection{Strengths}
Our framework has five main strengths.
\begin{itemize}
  \item It is system-driven: ecology, risk, access, and operations are linked explicitly.
  \item It accepts partial coverage as a design principle, which is realistic for large parks.
  \item It combines human and technological protection rather than treating them separately.
  \item It produces interpretable management outputs: hotspot maps, route plans, staffing needs, and sensitivity curves.
  \item It is transferable because the architecture is modular.
\end{itemize}

\subsection{Limitations}
The model also has limitations.
\begin{itemize}
  \item We use a stylized grid rather than full GIS raster and patrol track data.
  \item The poaching prior is based on exposure proxies because public incident records are incomplete.
  \item Ranger fatigue, corruption, and adversarial learning are represented only indirectly.
  \item Drone weather dependence and communication failures are not modeled in detail.
\end{itemize}

These limitations do not invalidate the framework, but they do define the next calibration steps for operational deployment.

\subsection{Most Critical Assumption}
The most sensitive assumption is the calibration of zone-level threat priors. If intelligence data are sparse or biased, hotspot ranking can shift. For this reason, we recommend periodic re-estimation of \(P_i\) using patrol logs, incident records, and station interviews before each planning cycle.

